%% PeerJ back matter. PeerJ collects the funding statement, competing interests,
%% ethics/consent, and author contributions in its submission form and asks that
%% they not appear in the manuscript; their text is kept verbatim in
%% "Peerj Submision/portal_statements.txt". Only the data and code availability
%% statements stay here, because they cite the repository and the paper's tables.

\section*{Data availability}

The benchmark artifacts reported in this paper are archived on Zenodo under
\url{https://doi.org/10.5281/zenodo.23765981} (release v1.1.0) \citep{Zenodo_RagRetrievalBench}
and maintained in the project repository \citep{GitHub_RagRetrievalBench}. The released artifacts comprise the source documents,
chunked corpora, generated synthetic queries, LLM-judged qrels including explicit negative
labels, per-model ranking files, raw-result JSON, and the generated dataset-level and
aggregate reports for all three corpora. Persisted ChromaDB vector indexes are excluded from
the release and can be rebuilt from the chunked corpora with the pipeline's indexing stage.
The source documents are public institutional regulations obtained from the providers listed
in Table~\ref{tbl:dataset_provenance} \citep{TED_Regulations, TED_Regulations_TR, BDDK_Mevzuat}; the
reproducibility conditions under which the reported snapshot was produced are summarized in
Table~\ref{tbl:reproducibility}.

\section*{Code availability}

The pipeline implementation, together with the configuration and stage sequence used to
produce the reported snapshot, is archived on Zenodo under
\url{https://doi.org/10.5281/zenodo.23765981} (release v1.1.0) \citep{Zenodo_RagRetrievalBench}
and maintained in the GitHub repository \url{https://github.com/kmkarakaya/RagRetrievalBench}
\citep{GitHub_RagRetrievalBench}. The repository README describes the code, the datasets,
the requirements, and the steps needed to reproduce the snapshot.
